% ======================================================================
% Personal essay — old newspaper style, two columns
% Adapted from the latex-paper-writing skill's essay template.
%
% IMPORTANT — engine: this file uses XeLaTeX, not pdfLaTeX.
% Reason: the essay text is Vietnamese, and Vietnamese precomposed
% diacritics (ệ, ố, ử, ẵ, ...) are NOT reliably covered by T1-encoded
% fonts (Palatino/mathpazo, IM Fell English) under pdfLaTeX. XeLaTeX +
% fontspec renders Unicode directly against an OpenType/TrueType font
% with full Vietnamese glyph coverage instead.
%
% Compile:  xelatex thien-tang.tex   (twice, for stable layout/rules)
%
% NOT compiled/verified in this environment (no TeX distribution
% installed here) — please compile and visually check before using.
% If a font below is missing, fontspec will error at \setmainfont;
% install one of Linux Libertine O / EB Garamond / Noto Serif, or let
% it fall through to Latin Modern Roman.
% ======================================================================
\documentclass[11pt,a4paper]{article}

\usepackage{fontspec}
\usepackage{polyglossia}
\setmainlanguage{vietnamese}

\usepackage[margin=1.6cm,top=1.4cm,bottom=1.8cm]{geometry}
\usepackage{multicol}
\usepackage{lettrine}        % drop caps — https://ctan.org/pkg/lettrine
\usepackage{microtype}       % protrusion/expansion: rescues narrow columns
\usepackage{setspace}
\usepackage{xcolor}
\definecolor{ink}{RGB}{26,22,18}   % slightly warm black, like aged ink
\color{ink}

% ---- Typeface: prefer an old-print serif with full Vietnamese coverage ----
\IfFontExistsTF{Linux Libertine O}{\setmainfont{Linux Libertine O}}{%
  \IfFontExistsTF{EB Garamond}{\setmainfont{EB Garamond}}{%
    \IfFontExistsTF{Noto Serif}{\setmainfont{Noto Serif}}{%
      \setmainfont{Latin Modern Roman}}}}

% ---- Justification tuning for narrow columns ----
\setlength{\columnsep}{18pt}
\setlength{\columnseprule}{0.4pt}   % thin vertical rule between columns
\emergencystretch=1.2em
\tolerance=2000
\hyphenpenalty=30

\setlength{\parindent}{1em}
\setlength{\parskip}{0pt}

% ---- Nameplate / masthead macros ----
\newcommand{\PaperName}{TRANG CHIỀU}
\newcommand{\PaperMotto}{``Ghi lại những trang sách đã đọc.''}
\newcommand{\PaperPlace}{HÀ NỘI}
\newcommand{\PaperDate}{THỨ BẢY, 11 THÁNG 7, 2026}
\newcommand{\PaperIssue}{Tập I \textperiodcentered{} Số 1}

\newcommand{\masthead}{%
  \begin{center}
    {\footnotesize\scshape \PaperIssue \hfill \PaperPlace \hfill \PaperDate}\\[2pt]
    \rule{\textwidth}{0.4pt}\\[4pt]
    {\fontsize{34}{36}\selectfont\scshape \PaperName}\\[3pt]
    {\small\itshape \PaperMotto}\\[4pt]
    \rule{\textwidth}{1.4pt}\\[-2.5pt]
    \rule{\textwidth}{0.4pt}
  \end{center}
}

% ---- Article headline + byline ----
\newcommand{\headline}[1]{%
  \begin{center}{\LARGE\scshape #1}\\[2pt]\rule{0.35\textwidth}{0.4pt}\end{center}}
\newcommand{\byline}[1]{%
  \begin{center}{\footnotesize\scshape Bởi #1}\\[2pt]\rule{0.2\textwidth}{0.4pt}\end{center}\vspace{4pt}}

% ---- Section break ornament (asterism, as in old print) ----
\newcommand{\sectionbreak}{%
  \par\vspace{8pt}{\centering *\quad*\quad*\par}\vspace{8pt}}

% ---- End-of-article rule ----
\newcommand{\closearticle}{%
  \par\vspace{6pt}{\centering\rule{0.5\columnwidth}{0.4pt}\par}\vspace{6pt}}

\pagestyle{empty}

\begin{document}

\masthead

\headline{Thiên Táng: Ba Mươi Năm Như Một Giấc Mộng Dài}
\byline{Thanh}

\begin{multicols}{2}

\lettrine[lines=3,lhang=0.05]{M}{ột quyển sách} mình đã bỏ dở khi đọc
được 30 trang hồi Tết năm 2026, đơn giản vì ở lúc đó mình không quá
hứng thú với câu chuyện mà nó truyền tải, cũng như ở thời điểm đó mình
có nhiều thú vui khác hơn, và có lẽ cũng do một phần mình đã đọc quá
nhiều, nên đã hình thành một cảm giác muốn tránh né việc đọc chăng?
Tuy nhiên dù lý do là gì, mình cũng đã khá lâu không tiếp tục đọc
quyển này. Khi tủ sách trong phòng mình bắt đầu dày lên, mình phải
tìm cách gì đó để giải quyết những quyển còn tồn đọng, và Thiên Táng
hiện lên như một lựa chọn khả dĩ, vì nó mỏng, và một phần mình cũng
muốn đổi gió khi đọc về những thứ mình chưa biết, về câu chuyện của
một tác giả nữ.

Không như lần đầu tiên, ở lần này mình đã thực sự đắm chìm trong câu
chuyện về Tây Tạng ở những năm giữa thế kỷ 20, khi mà xung đột giữa
họ và Trung Quốc được đẩy đến đỉnh điểm. Từng câu chữ của Hân Nhiên
như đưa mình qua từng tấc đất của phía tây Trung Quốc, nơi mình được
biết về văn hoá, phong tục, lối sống cũng như đức tin của người dân
Tây Tạng. Và mình nhận ra rằng vốn hiểu biết của mình nông cạn đến
như nào, và cũng như mình thực sự bị choáng ngợp trước chiều sâu của
một nền văn minh mà nếu xét theo tiêu chuẩn hiện tại, chúng ta đều
coi họ là ``lạc hậu''.

Trong bài này, mình sẽ không đi sâu vào chính trị Tây Tạng, tại sao
họ lại xung đột với Trung Quốc, cũng như lịch sử cận đại của họ.
Mình muốn dành thời gian để giới thiệu đến người đọc những hiểu biết
hạn hẹp của mình sau khi tìm hiểu về văn hoá cũng như phong tục của
Tây Tạng. Phần cuối mình sẽ dành để nói về bản thân tác phẩm, thứ mà
mình tin mọi người nên đọc thử một lần để thấy được vẻ đẹp của lời
văn của tác giả, cũng như vẻ đẹp của chính con người.

\sectionbreak

Cao nguyên Tây Tạng nằm ở phía tây Trung Hoa đại lục, nơi có độ cao
trung bình 4.500m, để dễ so sánh thì đỉnh Fansipan của Việt Nam cao
3.143m, thấp hơn cao nguyên này đến hơn 1km. Một người Tây Tạng từ
khi sinh ra đã sống ở một môi trường mà những người như chúng ta phải
dùng bình khí oxy để có thể sinh hoạt bình thường được. Ở độ cao này,
nồng độ oxy rất loãng, áp suất không khí cao, từ việc hít thở đến
nhóm lửa nấu ăn đều gặp rất nhiều khó khăn. Người Tây Tạng theo thời
gian đã tiến hoá để có thể thích nghi với những điều kiện khắc nghiệt
này.

Thời tiết trong một ngày cũng có thể biến động một cách mãnh liệt. Ở
cao nguyên trung tâm, nhiệt độ có thể lên tới 22°C, và nên nhớ rằng
cao nguyên Tây Tạng nằm gần bầu trời, nên tia UV cực mạnh và có tác
động lớn đối với họ. Ban đêm nhiệt độ có thể xuống tới $-8$°C, với
các cơn gió lạnh buốt da thịt. Ngoài ra Tây Tạng là một vùng đất khô,
có thể coi là sa mạc ở những vùng cao nguyên trung tâm, với lượng mưa
hàng năm rất thấp, chỉ từ 100mm cho tới 300mm. Từ đó không khí rất
khô, và nước uống cũng là vấn đề nan giải đối với con người. Mỗi ngày
với người Tây Tạng là một cuộc chiến chống lại thiên nhiên, khí hậu
để tìm cách sinh tồn trên mảnh đất quê hương của mình, và điều này
cũng sẽ ảnh hưởng trực tiếp đến tôn giáo, văn hoá cũng như phong tục
của họ, điều mà mình muốn bàn luận nhiều hơn ở các phần tiếp theo.

\sectionbreak

Tuỳ theo vị trí địa lý, mà người dân Tây Tạng sẽ có lối sống tương
ứng. Ở những vùng thung lũng, nơi điều kiện thích hợp hơn cho việc
trồng trọt thì người dân sẽ chọn lối sống định canh định cư, mình sẽ
không đề cập sâu hơn về nhóm này. Nhóm còn lại sống ở những thảo
nguyên cao, sống theo lối sống du mục, và di chuyển theo mùa với đàn
\emph{Yak}.

Tháng 3--5 hàng năm, sau một mùa đông khắc nghiệt thì đàn gia súc đã
rất yếu, trong khi đó thì cỏ chưa mọc kịp. \emph{Yak} bắt đầu sinh
con ở khoảng thời gian này, và nếu không ứng phó trước, thì nguồn thu
nhập của một gia đình có thể bay biến trong một đêm. Khi cỏ xuân mọc
lên vào cuối tháng 4, cuộc sống của dân du mục cũng trở nên dễ chịu
hơn. Họ di chuyển lên các bãi cỏ ở cao hơn, tiện cho việc chăn thả
gia súc và sản xuất bơ cũng như phô mai khô.

\vspace{6pt}
\begin{center}\footnotesize\itshape (Còn tiếp — kỳ sau.)\end{center}

\end{multicols}

\vfill
\begin{center}\rule{\textwidth}{0.4pt}\\[2pt]
{\footnotesize\scshape \PaperName\ \textperiodcentered{} in cho tác giả \textperiodcentered{} \PaperDate}
\end{center}

\end{document}
